\documentclass[11pt]{article}
\usepackage{fontspec}
\setmainfont{Times New Roman}
\setsansfont{Arial}
\setmonofont{Consolas}
\usepackage{amsmath,amssymb,mathtools}
\usepackage{geometry}
\usepackage{microtype}
\geometry{a4paper,margin=2.35cm}
\setlength{\parindent}{1.25em}
\setlength{\parskip}{0pt}
\makeatletter
\renewcommand\footnotesize{\@setfontsize\footnotesize{7.7}{8.7}\selectfont}
\makeatother
\emergencystretch=95em
\allowdisplaybreaks
\sloppy
\hbadness=10000
\vbadness=10000
\hfuzz=2pt
\vfuzz=10pt
\raggedbottom

\begin{document}

% Унит: Папер31 Сецтион03 ентры 4 в001. Дирецт Интерславиц Латин транслатион фром the Герман финал-аудитед соурце слице, соурце линес 202--220 / сегментс P31-С0046--P31-С0049. Фоотноте цоунтер цонтинуес афтер Папер31 Сецтион03 ентры 3.
\setcounter{footnote}{22}

2ц. Ако $\mathfrak R$ јест полно разложимо другого рода, тогда $\overline{\mathfrak R}$ не јест полно разложимо.

Како под 2б, достаточно јест предполагати, же $\mathfrak R$ јест полье, имено конечно разширјено полье другого рода над $P$. Треба показати, же в представјеньу $\overline{\mathfrak R}$ како директној сумы примарных колц наступа најманье једна властно примарна компонента. За то достаточно јест показати ненуловы елемент из $\overline{\mathfrak R}$, од кторого некоја степень изчезаје. Бо ако $\bar\beta\ne0$, тогда најманье једна компонента $\bar\beta_i\ne0$; ако $\bar\beta^\rho=0$, тогда $\bar\beta_i^\rho=0$ (\S 2, 2, хомоморфизм од $\overline{\mathfrak R}$ к $\overline{\mathfrak R}_i$), и зато $\overline{\mathfrak R}_i$ не јест полье.

Нехај $\gamma$ буде елемент другого рода из $\mathfrak R$, експоненты $f\ge 1$; нехај
\[
        F(\gamma)=\gamma^{kp^f}+c_1\gamma^{(k-1)p^f}+\cdots+c_k=0
\]
буде уравненье најнижшего степена, кторому $\gamma$ задовольа взходно к $P$, где $p$ означаје карактеристику $P$. Елементы
\[
        e,\gamma,\gamma^2,\ldots,\gamma^{kp^f-1}
\]
из $\mathfrak R$ сут зато линеарно независне взходно к $P$ и последовательно по дефиницији равности оставајут линеарно независне взходно к $\overline P$ такоже в $\overline{\mathfrak R}$. Ако зато положити
\[
        G(\gamma)=\gamma^{kp^{f-1}}+\sqrt[p]{c_1}\,\gamma^{(k-1)p^{f-1}}
        +\cdots+\sqrt[p]{c_k},
\]
тогда $G(\gamma)$ јест ненуловы елемент из $\overline{\mathfrak R}$; бо заради $p>1$ и $f\ge1$ число $kp^{f-1}$ јест всегда менше неж $kp^f$, зато не може быти никакој зависимости меджу степеньами $\gamma$, кторе наступајут в $G(\gamma)$. Але $G(\gamma)^p=F(\gamma)$ и зато изчезаје; с тым јест 2ц доказано.

Из 2а и 2ц следује, же полна разложимост првого рода $\overline{\mathfrak R}$ влече за собоју полну разложимост првого рода $\mathfrak R$; из 2б следује обрат. С тым јест теорем доказаны.

\end{document}

